% Blocchi migrati dal Cap. 4 per il vincolo di lunghezza del corpo.
% La rilevanza crittografica in forma estesa e' nell'Appendice C (Contesto).

Questa sezione raccoglie il materiale di contesto e di approfondimento dell'algoritmo di Shor: lo stato delle realizzazioni sperimentali, i miglioramenti algoritmici recenti, la roadmap verso il calcolo fault-tolerant e il confronto con l'algoritmo di Grover. Il Capitolo~\ref{chap:shor} ne riporta la sintesi.

\subsection{Sviluppi Recenti e Stato dell'Arte}

\subsubsection{Implementazioni su Hardware Reale}

Nonostante la solidità teorica, l'esecuzione di Shor su hardware reale rimane limitata a istanze molto piccole a causa dei requisiti di profondità circuitale. I principali risultati sperimentali sono:

\begin{itemize}
    \item \textbf{2001 -- \ac{NMR} (Vandersypen et al.)}: prima dimostrazione sperimentale di Shor su un sistema a 7 spin nucleari; fattorizzazione di $N = 15$ in $3 \times 5$ \cite{vandersypen2001}.
    \item \textbf{2016 -- Ioni intrappolati (Monz et al.)}: implementazione scalabile su 5 qubit ionici; fattorizzazione di $N = 15$ con un'architettura modulare \cite{monz2016}.
    \item \textbf{2021 -- Qubit superconduttori \ac{IBM}}: dimostrazione compilata della fattorizzazione di $N = 21$ su processori \ac{IBM} Quantum utilizzando 5 qubit \cite{skosana2021}.
\end{itemize}

Queste istanze evidenziano quanto la distanza dalle chiavi crittograficamente rilevanti sia ancora enorme. Una stima classica del fabbisogno, formulata sul modello di circuito privo di correzione d'errore, riporta $5K+1$ qubit di memoria e circa $72K^3$ gate elementari per fattorizzare un numero di $K$ bit \cite{beckman1996}; costanti e overhead dipendono però dall'architettura aritmetica e, nelle stime moderne, dal protocollo fault-tolerant, rispetto al quale quel conteggio va inteso come limite inferiore ottimistico. In ogni caso, i valori per RSA-2048 restano fuori portata per i processori \ac{NISQ} attuali.

Alla distanza teorica se ne aggiunge una operativa, documentata da una campagna recente su più piattaforme cloud: le costruzioni circuitali restano specifiche per ciascun modulo, e le fedeltà delle macchine si presentano instabili, con tassi d'errore elevati e fluttuanti \cite{shor_challenges2025}. È la ragione per cui, in questo lavoro, $N=15$ è trattato come istanza illustrativa e non come misura di fattibilità.

\subsubsection{Miglioramenti Algoritmici: Regev 2023}

Un risultato teorico rilevante è stato pubblicato nel 2023 da Oded Regev \cite{regev2023}. La proposta fattorizza un intero di $n$ bit eseguendo indipendentemente circa $\sqrt n+4$ circuiti, ciascuno con $\widetilde O(n^{3/2})$ gate, seguiti da post-processing classico polinomiale. Il confronto con Shor non può quindi essere ridotto al solo conteggio di gate di una singola esecuzione: aumentano il numero di run e, nella costruzione originale, anche lo spazio quantistico; inoltre l'articolo dichiara non chiaro l'effettivo vantaggio implementativo. Il risultato è un miglioramento asintotico importante, non una dimostrazione di praticabilità su dispositivi NISQ.

\subsubsection{Roadmap IBM verso il Quantum Fault-Tolerant}

\ac{IBM} ha delineato una roadmap concreta verso il calcolo quantistico fault-tolerant \cite{ibm_roadmap2025}:
\begin{itemize}
    \item \textbf{2025}: processori Loon e Nighthawk (120 qubit con 218 coupler sintonizzabili), con dimostrazione dei componenti chiave per la fault-tolerance tramite codici qLDPC;
    \item \textbf{2026}: target di vantaggio quantistico su problemi pratici;
    \item \textbf{2029}: sistema Quantum Starling con 200 qubit logici e oltre $10^8$ operazioni quantistiche.
\end{itemize}

Questi traguardi sono obiettivi industriali dichiarati da IBM, soggetti a revisione, e non stabiliscono quando una macchina avrà le risorse logiche necessarie a rompere RSA-2048. Rafforzano la necessità di preparare la migrazione, ma non forniscono una previsione crittografica.

Le sfide hardware principali per l'esecuzione di Shor su istanze crittograficamente rilevanti sono analizzate in dettaglio nel Capitolo~\ref{chap:rumore}, mentre l'implementazione pratica con Qiskit e le tecniche di ottimizzazione basate su intelligenza artificiale sono presentate nel Capitolo~\ref{chap:sviluppo}.
\clearpage

\subsection{Contesto: Confronto con l'Algoritmo di Grover}

Per collocare correttamente l'algoritmo di Shor nel panorama dell'informatica quantistica, è utile confrontarlo con l'altro grande risultato algoritmico degli anni Novanta: l'algoritmo di Grover per la ricerca non strutturata \cite{grover1996}.

\subsubsection{Il Problema di Grover}

Dato un insieme di $N$ elementi e una funzione oracolo $f:\{0,\ldots,N-1\}\to\{0,1\}$, l'obiettivo è trovare $x^*$ tale che $f(x^*)=1$. Classicamente, nel caso peggiore, sono necessarie $O(N)$ interrogazioni. Grover (1996) dimostrò che un algoritmo quantistico, sfruttando il meccanismo di \textit{amplitude amplification}, può trovare la soluzione in sole $O(\sqrt{N})$ interrogazioni -- uno \textbf{speedup quadratico} che è stato dimostrato ottimale: nessun algoritmo quantistico può fare meglio \cite{bennett1997,zalka1999}.

\subsubsection{Differenze Fondamentali}

I due algoritmi rappresentano due regimi distinti di vantaggio quantistico:

\begin{table}[H]
\centering
\small
\begin{tabular}{@{}lp{5cm}p{5cm}@{}}
\toprule
\textbf{Caratteristica} & \textbf{Shor} & \textbf{Grover} \\
\midrule
Problema target          & Fattorizzazione intera / period finding & Ricerca non strutturata \\[2pt]
Speedup                  & Esponenziale rispetto al miglior classico & Quadratico: $O(\sqrt{N})$ vs $O(N)$ \\[2pt]
Tecnica quantistica core & QFT + Phase Estimation          & Amplitude Amplification \\[2pt]
Ottimalità               & Non dimostrata in assoluto      & Dimostrata ottimale \cite{bennett1997} \\[2pt]
Impatto crittografico    & Rompe RSA, DH, ECDSA (chiave pubblica) & Dimezza sicurezza AES/SHA (simmetrica) \\[2pt]
Requisiti hardware       & $O(\log N)$ qubit, alta profondità & $O(\log N)$ qubit, bassa profondità \\
\bottomrule
\end{tabular}
\caption{Confronto tra l'algoritmo di Shor e l'algoritmo di Grover}
\end{table}

L'impatto sulla crittografia simmetrica di Grover è significativo ma gestibile: raddoppiare la lunghezza delle chiavi (da AES-128 a AES-256) è sufficiente a ripristinare il livello di sicurezza. Al contrario, non esiste un analogo rimedio per RSA: l'unica risposta è la migrazione verso gli standard post-quantum discussi nella Sezione~\ref{app:intro:pqc}. È per questa ragione che l'algoritmo di Shor rappresenta la minaccia quantistica più urgente e motivante per la ricerca sul calcolo quantistico applicato.

Tutta l'analisi svolta in questa sezione ha però assunto un modello di calcolo \textit{ideale}: porte unitarie perfette, coerenza illimitata, nessuna interazione con l'ambiente. Su hardware reale questi presupposti non reggono, e la distanza tra il modello teorico e l'esecuzione fisica è precisamente il nodo critico che rende ancora irraggiungibile la fattorizzazione di numeri crittograficamente rilevanti. Comprendere le origini e gli effetti del rumore quantistico è quindi il passo necessario per valutare le reali prospettive dell'algoritmo di Shor.

\subsection{Sensibilità delle componenti di Shor al rumore}
\label{app:shor_rumore_esteso}

La seguente analisi è stata trasferita dalla Sezione sul rumore del
Capitolo~\ref{chap:quadro}. Nel corpo resta il nesso causale essenziale; qui si
conservano derivazioni, esempio numerico e figura di scala.

\subsubsection{Errori nella Quantum Fourier Transform}

La \ac{QFT} usa rotazioni
$R_k=\operatorname{diag}(1,e^{2\pi i/2^k})$ con angoli esponenzialmente
piccoli. Per $k=10$, $2\pi/1024\simeq6\times10^{-3}$ rad. Gli errori coerenti
di calibrazione possono sommarsi in ampiezza, mentre il canale depolarizzante
produce degradazione stocastica: non è quindi corretto ricondurli a un solo
``errore di fase'' lineare. Le campagne li mantengono distinti.

\subsubsection{Errori nella Phase Estimation}

I qubit di controllo devono mantenere la coerenza lungo il cammino critico
dell'esponenziazione modulare. Come ordine di grandezza,
\begin{equation}
    T_2\gg\tau_{\text{gate}}\,O(n^3),
\end{equation}
ma questa espressione non determina una soglia universale in bit: la costante
dipende da aritmetica e compilazione, mentre in regime fault-tolerant il
confronto diretto è sostituito da un budget di errore logico.

\subsubsection{Errori nell'esponenziazione modulare}

Sotto l'ipotesi puramente illustrativa di $L$ eventi indipendenti con
probabilità $\epsilon$, il proxy di nessun evento è
\begin{equation}
    P_{0\text{-eventi}}=(1-\epsilon)^L,\qquad L\sim O(n^3).
\end{equation}
Ponendo soltanto per visualizzare la scala $L=n^3$, per $n=15$ e
$\epsilon=10^{-3}$ si ottiene
$(1-10^{-3})^{3375}\simeq0.034$. Non è una probabilità di successo di Shor:
alcuni errori possono non alterare l'esito e la formula non rappresenta
rilassamento, readout o struttura circuitale.

\begin{figure}[H]
\centering
\includegraphics[width=0.8\textwidth]{15_psuccess_vs_n.pdf}
\caption[Proxy di nessun evento al crescere del circuito]{Proxy illustrativo
$(1-\epsilon)^{n^3}$ per cinque valori di $\epsilon$. La curva non deriva dai
circuiti sperimentali e mostra soltanto l'accumulo moltiplicativo di eventi.}
\label{fig:psuccess_vs_n}
\end{figure}

\subsection{Scelta dell'aritmetica e perimetro delle istanze}
\label{app:shor_perimetro_esteso}

Questo dettaglio metodologico proviene dalla sezione sul disegno sperimentale.
Sintetizzare l'esponenziazione modulare come un operatore unitario denso costa
$O(4^k)$ porte elementari su $k$ qubit~\cite{Barenco1995}. La decomposizione
di Beauregard~\cite{beauregard2002} evita tale sintesi costruendo l'aritmetica
nel dominio della \ac{QFT}, con addizionatori di Draper, e richiede
$O(n^3)$ porte su $2n+3$ qubit. Il costo resta elevato: per $N=21$, $a=2$ e
otto qubit di conteggio, la compilazione congiunta in
\texttt{rz/sx/x/cx} produce 21\,036 porte \ac{CX} e profondità 23\,081.
Il costo dominante non è la trasformata finale della \ac{QPE}, ma le
trasformate interne agli addizionatori; è lì che l'approssimazione deve agire
per incidere sul conteggio di porte.

Le campagne rumorose sono perciò limitate a $N=15$, $a=7$. Le istanze
$N=21$, $a=2$, e $N=35$, $a=6$, sono state validate idealmente mediante truth
table, pulizia delle ancille e distribuzione QPE, ma escluse dal campionamento
rumoroso per il costo della simulazione. Questa distinzione impedisce di
confondere validazione funzionale e fattibilità rumorosa.

In termini di spazio, fattorizzare un intero di $n$ bit richiede almeno
$2n$ qubit logici per la phase estimation, oltre alle ancille necessarie
all'aritmetica modulare. Il conteggio non include l'overhead fault-tolerant e
non basta a stimare la fattibilità: connettività e routing possono aumentare
profondità e numero di porte senza modificare il numero di qubit logici.
